% !TeX root = ../../../main.tex
기룡이는 친구들과 함께 경기대학교 정문 쪽에 위치한 컨벤션 센터를 대관하여 대탈출 게임을 하려고 한다.
대탈출 게임의 규칙은 아래와 같다.

게임은 1명의 심판과 나머지 탈출자들로 구성된다.
탈출자는 몇 명이든 상관없지만 심판 역할은 반드시 1명만 배정되어야 한다.
역할 분배가 완료되면 탈출자들은 잠시 밖에서 대기한다.

심판은 양의 정수 $n$을 선택하여, 좌석이 $n \times n$개 배열된 영역을 노란색 테이프로 표시한다. 이 영역을 `탈출 영역'이라고 부른다.
이후 탈출 영역 안의 원하는 좌석에 `X' 표시를 한다.
탈출 영역의 준비가 끝나면 탈출자들은 순서대로 한 명씩 컨벤션 센터에 들어와서 탈출을 시도한다.

탈출 영역의 왼쪽 아래 모서리의 좌석에서 시작하여 오른쪽 위 모서리의 좌석에 도달하면 탈출에 성공한다.
심판은 각각의 탈출자들이 탈출에 성공하기까지 몇 번 이동했는지 기록한다.
한 좌석에서 다른 좌석으로 위치를 옮겼을 때 1회 이동한 것으로 본다.
현재 좌석에서 상하좌우로 인접한 좌석으로만 이동할 수 있다.
그러나 `X' 표시가 된 좌석으로는 이동할 수 없고, 탈출 영역을 벗어나서는 안 된다.

이때 탈출하기까지 최대한 적은 횟수로 이동해야 더 높은 등수를 얻을 수 있으므로, 가장 적은 횟수로 이동한 탈출자가 우승하게 된다.

탈출자가 된 기룡이는 이번 게임에서 너무 1등을 하고 싶은 나머지, 탈출 영역을 한눈에 볼 수 있도록 목을 길게 빼는 기술을 연마했다.
기룡이가 탈출할 차례가 되었을 때, 기룡이는 목을 하늘 높이 빼서 탈출 영역을 한눈에 파악하였다.

기룡이가 파악한 탈출 영역이 2차원 배열로 주어졌을 때, 탈출에 필요한 최소 이동 횟수를 계산해 보자.
